\cvsection{Formation Académique}

\cveducation{2025-2027 \\ (en cours)}{Université de Sherbrooke}{Sherbrooke}
{Maîtrise de type recherche}{Génie électrique}

\cveducation{2020-2024}{Université de Sherbrooke}{Sherbrooke}
{Diplôme d'études universitaires}{Génie informatique}

\cveducation{2018-2020}{Champlain College}{Saint-Lambert}
{Diplôme d'études collégiales}{Computer Sciences and Mathematics}

\cvsection{Sujet de recherche}

\cvexperience{2025-2027 \\ (en cours)}{Convertisseur analogique-numérique sur logique programmable pour la lecture dispersive de qubits supraconducteurs}
{Institut Quantique de l'Université de Sherbrooke}{Sherbrooke}
{Conception d'un ADC sur FPGA pour la lecture de signaux analogiques en milieu cryogénique}
{\vspace{-12pt}
\begin{itemize}
    \item Intégration de la \textit{toolchain} pour la programmation et la vérification du FPGA en cryogénie :
    \begin{itemize}
        \item \textbf{LiberoSOC} : Outil de synthèse, placement, routage et programmation
        \item \textbf{QuestaSIM, cocotb} : Outil et librairie permettant la simulation \textit{HDL}
        \item \textbf{Tessent} : Suite d'outils d'insertion de chaîne de tests \textit{IJTAG}
        \item \textbf{PyHegel} : Bibliothèque Python de contrôle des instruments de laboratoire pour les essais
    \end{itemize}
    \item Conception de \textit{PCBs} mixtes analogique-numérique pour l'acquisition de signaux du cryostat
    \item Formation et supervision de stagiaires au CEGEP et à l'université
    % \begin{itemize}
    %     \item Filtrage de signaux provenant d'expérience en laboratoire en Python
    %     \item Création de modèles pour LTSPICE de pièces utilisé dans le circuit avec MATLAB
    % \end{itemize}
    % \item Élaboration du cahier des charges avec une équipes de chercheur multidisciplinaire
    % \item Utilisation des équipements d'assemblage et de tests des \textit{PCBs}
    % \item Dessin de la schématique et du routage d'un PCB 8 couches haute vitesse
\end{itemize}}
